% \documentclass[12pt,largemargins]{homework}

% \begin{document}
    
% \question#{0}


% محاسبه و شمارش تعداد توالی‌هایی که ساختارهای خاص بیولوژیکی را تشکیل می‌دهند، برای ارزیابی احتمال وقوع جهش‌های موفق در ویروس‌ها اهمیت فراوانی دارد. به عنوان مثال، در مولکول‌های تک‌رشته‌ای $RNA$، برای ایجاد پایداری، رشته روی خودش تا می‌خورد و ساختاری به نام «سنجاق‌سر» (\lr{Hairpin} یا \lr{Stem-Loop}) تشکیل می‌دهد. 

% برای تشکیل این ساختار، یک بخش از رشته (ساقه اولیه) باید دقیقاً با بخش دیگری در ادامه رشته (ساقه ثانویه) جفت شود و بین این دو بخش، یک حلقه (Loop) قرار می‌گیرد. قانون جفت شدن بازها در $RNA$ به این صورت است که $A$ فقط با $U$، و $C$ فقط با $G$ جفت می‌شود. 
% بنابراین اگر ساقه اولیه را رشته $w$ بنامیم، ساقه ثانویه باید دقیقاً «معکوسِ مکملِ» آن باشد (آن را $w^{RC}$ می‌نامیم). 

% الفبای ما شامل چهار نوکلئوتید $RNA$ است: $\Sigma = \{A, C, G, U\}$.
% زبان تولیدکننده این ساختارهای سنجاق‌سری را می‌توان به شکل زیر مدل کرد (فرض کنید طول رشته حلقه وسط دقیقاً ۵ نوکلئوتید است):
% $$L = \{ w v w^{RC} \mid w \in \Sigma^* \,,\, v \in \Sigma^5 \}$$
% (به عنوان مثال اگر $w = ACG$ باشد، معکوس مکمل آن $w^{RC} = CGU$ خواهد بود).

% \begin{itemize}
%     \item[الف)] آیا زبان $L$ که این ساختارهای $RNA$ را تولید می‌کند، یک زبان منظم (\lr{Regular}) است؟ در این صورت دلیل خود را بیاورید.
    
%     \item[ب)] اگر طول ساقه $w$ برابر با $k$ نوکلئوتید باشد (در نتیجه طول کل رشته $n = 2k + 5$ خواهد بود)، به طور کلی چند توالی $RNA$ مختلف وجود دارد که بتواند این ساختار سنجاق‌سری معتبر را تشکیل دهد؟
    
%     \item[ج)] فرض کنید یک ویروس برای فرار از دست سیستم ایمنی میزبان، به طور تصادفی رشته‌های $RNA$ با طول کل $n = 2k + 5$ تولید می‌کند تا حداقل یکی از آن‌ها ساختار سنجاق‌سری معتبر (زبان $L$) را تشکیل دهد. سیستم ایمنی میزبان در برابر این جهش‌ها مکانیزم دفاعی زیر را دارد:
%     به ازای تولید هر ۱۰ دسته از $RNA$های نامعتبر و بی‌اثر توسط ویروس، سیستم ایمنی هوشیارتر شده و زمان لازم برای تولید دسته‌های بعدی ویروس دو برابر می‌شود (۱ روز، سپس ۲ روز، سپس ۴ روز و الی آخر). 
%     با فرض اینکه ویروس بتواند حداکثر ۱۰ سال در بدن زنده بماند، حداقل طول ساقه ($k$) چقدر باید باشد تا احتمال اینکه جهش‌های تصادفی ویروس بتوانند حتی یک ساختار معتبر بسازند، به زیر $0/001$ درصد برسد؟(
% برای مقادیر بسیار کوچک $P$، از تقریب $1 - (1 - P)^M \approx M \times P$ استفاده می‌کنیم)
% \end{itemize}

% \ifcompileanswers\solution
% \textbf{الف)} 
% خیر، زبان $L$ منظم نیست. برای بررسی اینکه رشته انتهایی دقیقاً معکوس مکمل رشته ابتدایی ($w^{RC}$) هست یا نه، اتوماتا باید کل زیررشته $w$ (که طول آن از پیش تعیین نشده و می‌تواند هر مقداری باشد) را به خاطر بسپارد. اتوماتاهای متناهی حافظه محدودی دارند و نمی‌توانند وابستگی‌های با طول نامتناهی را مدیریت کنند. این موضوع را می‌توان با استفاده از لم تزریق (\lr{Pumping Lemma}) برای زبان‌های منظم نیز به طور دقیق اثبات کرد.

% \textbf{ب)}
% طول رشته ساقه اولیه ($w$) برابر $k$ است. از آنجا که الفبای ما ۴ حرف دارد ($\Sigma = \{A, C, G, U\}$)، تعداد $4^k$ حالت مختلف برای انتخاب رشته $w$ وجود دارد. به محض اینکه رشته $w$ انتخاب شود، رشته ساقه ثانویه ($w^{RC}$) به صورت کاملاً یکتا و قطعی مشخص می‌شود (یعنی تنها ۱ حالت برای آن وجود دارد).
% برای بخش حلقه ($v$) که طول آن دقیقاً ۵ است نیز $4^5$ حالت مختلف وجود دارد. از آنجا که انتخاب $w$ و $v$ مستقل از هم هستند، طبق اصل ضرب، کل توالی‌های معتبر برابر است با:
% $$ \text{Total Valid Sequences} = 4^k \times 4^5 \times 1 = 4^{k+5} $$

% \textbf{ج)}
% ابتدا باید ببینیم ویروس در مدت ۱۰ سال (حدود ۳۶۵۰ روز) حداکثر چند دسته رشته می‌تواند تولید کند.
% \begin{itemize}
%     \item ۱۰ دسته اول: هر کدام ۱ روز $\rightarrow$ مجموعاً ۱۰ روز
%     \item ۱۰ دسته دوم: هر کدام ۲ روز $\rightarrow$ مجموعاً ۲۰ روز
%     \item ۱۰ دسته سوم: هر کدام ۴ روز $\rightarrow$ مجموعاً ۴۰ روز
% \end{itemize}
% به طور کلی زمان لازم برای $m$ گروهِ ۱۰تایی از تولیدات برابر است با مجموع تصاعد هندسی:
% $$ 10 \times (1 + 2 + 4 + \dots + 2^{m-1}) = 10 \times (2^m - 1) $$
% می‌خواهیم این زمان از ۳۶۵۰ روز بیشتر نشود:
% $$ 10 \times (2^m - 1) \le 3650 \implies 2^m - 1 \le 365 \implies 2^m \le 366 $$
% جواب را به صورت تقریبی 9 در نظر می‌گیریم.
% اکنون احتمال موفقیت را حساب می‌کنیم. کل رشته‌های ممکن با طول $n = 2k + 5$ برابر $4^{2k+5}$ است. احتمال ($P$) اینکه یک رشته تولید شده تصادفی معتبر باشد برابر است با:
% $$ P = \frac{4^{k+5}}{4^{2k+5}} = \frac{1}{4^k} = 4^{-k} $$
% احتمال اینکه ویروس در $M=90$ تلاش خود حداقل یک رشته معتبر بسازد، باید کمتر از $0/001$ درصد ($10^{-5}$) باشد:
% $$ P_{\text{success}} = 1 - (1 - P)^{90} < 10^{-5} $$
% برای مقادیر بسیار کوچک $P$، از تقریب $1 - (1 - P)^M \approx M \times P$ استفاده می‌کنیم:
% $$ 90 \times 4^{-k} < 10^{-5} \implies 4^k > 90 \times 10^5 = 9 \times 10^6 $$
% همان‌طور که مشخص است، برای شرط $4^k > 9,000,000$، حداقل مقدار صحیح برای $k$ برابر با \textbf{۱۲} است. 
% بنابراین حداقل طول ساقه باید $k=12$ نوکلئوتید باشد.

% \fi\end{document}
